\paragraph{Bind a named method to the event that actually occurred.}
The declared book separates company-specific accounts from a method's financial
roles. Selecting that method also requires an unambiguous connection to the event.
\texttt{CompanyMethodRegistry/1} names the available data-only method definitions,
their schemas, expected account classifications and input types. Each entry
identifies its complete method body. The optional \texttt{CompanyMethodCall/1}
selects one entry under the book's company context and checks every entry's
structure before evaluating the selected method's prerequisites. An invalid unused
entry therefore cannot remain hidden until a later event happens to select it.

Some inputs are supplied facts; others describe the event itself. Keeping these
sources distinct matters for a sale after a tax close. Suppose the event is dated
15 June, but a separately supplied date says 1 September. A predicate comparing
that supplied date with a 1 July close could accept the wrong timing. The named
binder instead derives designated date and identity inputs from the event's own
metadata. It rejects a supplied input with the same name. In the checked case,
the 15 June sale consequently fails the closed-period prerequisite. A genuine
change to the event date changes the predicate's input with it.

The binder materializes the ordinary company-book transition request, including
the current state and selected method identities. The existing prerequisite,
monetary and memory engines then determine the result. Two company payment layouts
and a complete-lot sale reproduce their explicitly authored requests and entire
provisional transitions; unknown and false applicability findings retain their
different consequences. Repeated calls read the resulting book, so the next
payment uses the reduced outstanding amount. Application recomputes the complete
named call and uses the same type-sensitive JSON comparison as the underlying
transition.

The current registry supports a directly dated event. Assessment intervals,
measurement dates and period carry are different clock families and remain
unsupported here. The registry also leaves native applicability to an explicit
adapter. For example, a reusable inventory-sale adapter must validate all lot
allocations, disposition and tax-base conditions, the selected sale method and the
native close date, as well as preserve every unmapped field. Matching one native
journal cannot establish those conditions. The selected comparisons establish the
named binding mechanism; they do not turn the retained test projections into
general native-method admission.
